\scanpage{204}
Endomorphismes de l'anneau différentiel $\frac12$ simplicial $\mathcal A_*$

(NB. on ne demande plus de respecter la structure de $S$-algèbre, $S=\mathcal A_0$, i.e. on ne demande plus que $\varphi_*$ induise l'identité sur $\mathcal A_0$).

Rappelons-nous qu'un complexe de De Rham peut se construire \unsure{dès que l'on a} un anneau comm. $S$ avec idéal $S^+$ à puiss. divisées, et $t\in S^+$, i.e. une \deleted{\unsure{algèbre}} \inserted{$S$-algèbre} différentielle \inserted{$\mathbb N$-graduée, à p.d.} $\frac12$ simpliciale $\mathcal A_*(S,S^+,t)$, où
\[
\mathcal A_n(S,S^+,t)=S\{X_0,-,X_n\}[dX_0,-,dX_n]/\bigl(\sum X_i-t,\ \sum dX_i\bigr)_{pd}
\]
Si on a un hom. d'anneaux à p.d.
\[
S\xrightarrow{u}S'\quad(u(S^+)\subset S'^+-)\quad\text{et}\quad\lambda'\in S'
\]
tel que \deleted{$\unsure{\lambda\in S'^+}$} $u(t)=\lambda t'$, on définit un hom \inserted{de $S$-\unsure{alg. $\mathbb N$-graduées} différentielles à p.d.}
\[
\mathcal A_*(u,\lambda):\mathcal A_*(S,S^+,t)-\mathcal A_*(S',S'^+,t'),
\]
qui en degré $n$ est l'hom.
\diagram{\begin{tikzcd}
S\{X_0,-,X_n\}[dX_0,-,dX_n]/-\arrow[d]\\
S'\{X_0,-,X_n\}[dX_0,-,dX_n]/-
\end{tikzcd}}
donné par \deleted{$\unsure{X_0\longmapsto}$} $u:S-S'$ sur les scalaires,
\[
X_0\longmapsto\lambda X_0\quad\left(\text{donc }\sum X_0-t\longmapsto\sum\lambda X_0-u(t)=\sum\lambda X_0-\lambda t'=\lambda(\sum X_0-t')\right)
\]
$dX_0\longmapsto\lambda dX_0$. Ainsi $\mathcal A_*(S,S^+,t)$ devient un foncteur covariant de la catégorie des triples $(S,S^+,t)$, où les hom sont les couples $(u,\lambda)$ \inserted{\fbox{comme dessus}} (\unsure{composés} par $(u',\lambda')(u,\lambda)=(u'u,u'(\lambda)\lambda')$).

\uline{NB} En plus de sa graduation, $\mathcal A_n(S,S^+,t)$ admet une filtration \inserted{croissante} naturelle, où $\operatorname{Fil}^{\unsure{n}}\mathcal A_n(S,S^+,t)$ est la somme des expressions
\[
\sum F_{i_0- i_p}\,dX_{i_0}\wedge-\wedge dX_{i_p},
\]
avec $F_{i_0-i_p}\in S\{X_0,-,X_n\}$ étant de ``degré'' $\leq\unsure{n-r|n-p}$, ce sont des sous-$S$ modules (mais pas des idéaux), on a
\[
\operatorname{Fil}^n.\operatorname{Fil}^m\subset\operatorname{Fil}^{n+m},\quad\operatorname{Fil}^0=S\ (\text{scalaires}),\quad\operatorname{Fil}^n=0\text{ si }n<0.
\]
Cette filtration est respectée par les $\mathcal A_*(u,\lambda)$.
\margin{203}
